\Article{368}
Wir sagen, eine Kongruenz $S\equiv0$ habe \emph{gleiche} Wurzeln oder allgemeiner gleiche Teiler, wenn sie durch eine Potenz irgendeiner Funktion teilbar ist.

\begingroup\postdisplaypenalty=10000\SourceLeaf{246}
Ob eine gegebene Kongruenz gleiche Teiler hat, entscheidet man ebenso wie in der Theorie der Gleichungen. Setzen wir
\[
X\equiv\xi^mP
\]
so ist offensichtlich
\[
\frac{\mathrm{d}X}{\mathrm{d}x}\equiv\xi^{m-1}\left(mP\frac{\mathrm{d}\xi}{\mathrm{d}x}+\xi\frac{\mathrm{d}P}{\mathrm{d}x}\right)
\]
also ist $\frac{\mathrm{d}X}{\mathrm{d}x}$ durch $\xi^{m-1}$ teilbar. Allgemeiner sei
\[
X\equiv A^aB^bC^c\text{ usw.}
\]
wobei $A,B,C$ usw. verschiedene Primfunktionen bezeichnen. Dann ist
\[
\frac{\mathrm{d}X}{\mathrm{d}x}\equiv X\left(\frac{a\,\mathrm{d}A}{A\,\mathrm{d}x}+\frac{b\,\mathrm{d}B}{B\,\mathrm{d}x}+\frac{c\,\mathrm{d}C}{C\,\mathrm{d}x}+\text{usw.}\right)
\]
Daraus folgt, dass, sofern nicht eine der Zahlen $a,b,c$ usw. durch den Modul teilbar ist, $\frac{\mathrm{d}X}{\mathrm{d}x}$ durch $A^{a-1}B^{b-1}C^{c-1}$ usw., nicht aber durch $A^a,B^b,C^c$ usw. teilbar ist. Hieraus folgt der\par\endgroup

\textsc{Satz.} \emph{Ist der gemeinsame Teiler höchsten Grades der Funktionen $X$ und $\frac{\mathrm{d}X}{\mathrm{d}x}$ gleich $\xi$, so kommen alle Primfaktoren von $\xi$ auch in $X$ vor, und zwar jeweils mit einer um $+1$ größeren Vielfachheit als in $\xi$. Sind also $X$ und $\frac{\mathrm{d}X}{\mathrm{d}x}$ teilerfremde Funktionen, so hat $X$ keine gleichen Faktoren.}

\Article{369}
Beispiel I. Gesucht ist, ob die Funktion
\[
x^5+3x^4-6x^3+3x-4\ldots(X)
\]
modulo $17$ gleiche Teiler hat. Es ist
\[
\frac{\mathrm{d}X}{\mathrm{d}x}\equiv5x^4-5x^3-xx+3
\]
Hieraus findet man, dass die Funktionen $X$ und $\frac{\mathrm{d}X}{\mathrm{d}x}$ teilerfremd sind, sodass $X$ keine gleichen Teiler hat.

\Needspace{14\baselineskip}Beispiel II. Es sei
\[
X\equiv x^5+6x^4-3x^3-4xx+2x-3(\mathrm{mod.}\,13)
\]
Dann ist
\[
\frac{\mathrm{d}X}{\mathrm{d}x}\equiv5x^4-2x^3+4xx+5x+2
\]
Der größte gemeinsame Teiler der Funktionen $X,\frac{\mathrm{d}X}{\mathrm{d}x}$ ist $\equiv5xx+7x+7$ oder,
\SourceLeaf{247}
mit $8$ multipliziert: $xx+4x+4$. Da dieser Teiler $\equiv(x+2)^2$ ist, ist die Funktion $X$ durch $(x+2)^3$ teilbar, und der Quotient (der $xx+11$ ist) enthält keinen weiteren doppelten Teiler.

\Needspace{24\baselineskip}\Article{370.\ 371}
Ist die Funktion $X$ nach den vorigen Abschnitten als $A^aB^bC^c$ usw. dargestellt, wobei $A,B,C$ usw. teilerfremd und die Zahlen $a,b,c$ usw. ungleich sind, so lässt sich die Zerlegung weiter fortsetzen. Es sei also $X$ eine Funktion, die keine gleichen Teiler mehr enthält. Wir haben oben gesehen, dass $x^p-x$ das Produkt aller Primfunktionen ersten Grades ist. Es sei $\xi$ der gemeinsame Teiler höchsten Grades der Funktionen $X$ und $x^p-x$; dann ist $\xi$ das Produkt aller Teiler von $X$ ersten Grades, und $\frac X\xi$ hat keine solchen Teiler mehr. Findet man dagegen, dass die Funktionen $X$ und $x^p-x$ teilerfremd sind, so hat $X$ keinen Teiler ersten Grades, und daher hat die Kongruenz $X\equiv0$ keine wirklichen Wurzeln. Da ferner $x^{pp}-x$ das Produkt aller Primfunktionen zweiten und ersten Grades ist, enthält der gemeinsame Teiler höchsten Grades der Funktionen $x^{pp}-x$ und $\frac X\xi$, nämlich $\xi'$, alle Teiler von $X$ zweiten Grades. Setzt man dies fort, so erkennt man, dass $X$ auf diese Weise in Faktoren $\xi,\xi',\xi''$ usw. zerlegt wird, die jeweils alle Teiler ersten, zweiten, dritten usw. Grades enthalten.

\Article{372}
Ist aber das Produkt mehrerer Primfunktionen gleichen Grades gegeben, so müssen die einzelnen Funktionen durch Probieren gefunden werden. Diese Aufgabe hat große Ähnlichkeit mit derjenigen, die Faktoren zusammengesetzter Zahlen zu finden. Hier ist jedoch schon a priori bestimmt, ob sich die gegebene Funktion noch in Faktoren zerlegen lässt. Da auch hier die Anzahl aller möglichen Faktoren endlich ist, können wir ein ähnliches Hilfsmittel wie oben verwenden. Doch wollen wir uns damit nicht aufhalten, denn ein geübter Rechner, der die Prinzipien gründlich erfasst hat, wird bei Bedarf leicht besondere Kunstgriffe finden.

Wir gehen zu einem anderen Thema über, nämlich zur Betrachtung von Kongruenzen, deren Modul keine Primzahl ist, wie wir es bisher stets vorausgesetzt haben. Besonders verdient hier der Fall Aufmerksamkeit, dass der Modul eine Potenz einer Primzahl ist, sowohl an sich als auch deshalb, weil er zur Beseitigung einiger Zweifel (\textsection.\textsection.\ .\ .) nötig ist.
